\section{Capacity- and Compute-Matched Comparison with FedAvg}
\label{app:capacity_compute_comparison}

\subsection{Digits setting and model}

We use the five-domain Digits benchmark associated with
\citet{li2021fedbn}: MNIST, SVHN, USPS, SynthDigits and MNIST-M.
Each domain is assigned to one client with 743 training images, giving
3,715 training images in total. The fixed training partition and the full
domain test sets are used for both methods. Images are represented as
three-channel $28\times28$ inputs, with pixel values normalized to
$[-1,1]$. All five clients participate in each of 300 communication rounds,
and shared states are aggregated uniformly.

The classifier has three convolutional layers with channel widths
$3,64,64,128$, followed by fully connected layers with widths
$6272,2048,512,10$. Its batch-normalization parameters and buffers are shared.
FusedSpaceFed uses a compact U-Net with base width 16 and bottleneck
width $d_z=64$, a persistent private encoder for each client, and a shared
decoder. Classification uses the fused input $x+D(E_i(x))$.

For FedAvg, the first fully connected layer is widened from 2048 to 2065
units; the following layer and the corresponding batch-normalization
layer are adjusted accordingly. The active parameter count per client is
14,336,285 for FusedSpaceFed, comprising 14,219,210 classifier parameters,
72,080 encoder parameters and 44,995 decoder parameters.
The enlarged FedAvg classifier has 14,334,589 parameters, a relative
difference of $0.01183\%$. FusedSpaceFed stores its five client encoders
separately, together with the shared decoder and classifier.

\subsection{Local optimization and FLOP budget}

Each FusedSpaceFed client performs one encoder-only reconstruction epoch
with the decoder frozen, followed by one classification epoch updating
the encoder, decoder and classifier. Reconstruction uses mean squared
error; classification uses cross-entropy. The maximum batch size is 32,
and the final batch of each epoch is retained.

The FedAvg local workload is set using the combined reconstruction and
classification FLOP budget of FusedSpaceFed for each client and round.
It comprises one pass through the local training set followed by additional
shuffled minibatches. Batch sizes are selected using integer arithmetic,
and the unspent remainder is carried to subsequent rounds.

Dense forward and backward operations are counted with Torch 2.5.1
\texttt{FlopCounterMode}, assigning two FLOPs to a fused multiply-add.
The count includes gradient propagation through the frozen decoder during
reconstruction. Optimizer updates and gradient-norm clipping are accounted
for using $5P$ scalar operations per SGD step and $16P$ per Adam step,
where $P$ is the number of parameters updated by that optimizer.
Batch normalization, activations, pooling, loss evaluation, finite-value
checks, copies, aggregation and kernel overhead are outside this FLOP
accounting convention.

Table~\ref{tab:capacity_compute_resources} gives the cumulative counts
per final run. The relative difference in counted training FLOPs is
$0.000198\%$. Across all clients and rounds, FusedSpaceFed performs
36,000 reconstruction steps and 36,000 classification steps; FedAvg
performs 63,000 classification steps. The corresponding numbers of
supervised sample presentations are 1,114,500 and 1,956,535.

\begin{table}[H]
\coauthorcolor
\centering
\small
\setlength{\tabcolsep}{4pt}
\caption{Active parameters per client and cumulative training FLOPs
per run. The FusedSpaceFed FLOP count includes reconstruction warm-up.}
\label{tab:capacity_compute_resources}
\begin{tabular}{lrr}
\toprule
Method & Parameters & FLOPs ($10^{12}$) \\
\midrule
FusedSpaceFed & 14,336,285 & 551.898005 \\
FedAvg, enlarged & 14,334,589 & 551.896914 \\
\bottomrule
\end{tabular}
\end{table}

\subsection{Validation-based calibration}

For each domain, 149 of the 743 training images are assigned to
validation and 594 to fitting, using deterministic class-stratified
selection with data seed 20261005. The same split is used by both
methods.

For FedAvg, nine configurations are evaluated for 60 rounds with
seed 142, combining classifier learning rates
$\{0.005,0.01,0.02\}$ and gradient-clipping norms $\{0.5,1,2\}$.
Selection uses the uniform mean of domain validation accuracies at
round 60 and yields classifier learning rate 0.02 and clipping norm 2.

FusedSpaceFed uses the validation-selected configuration described in
Appendix~\ref{app:feature_heterogeneity_digits}. Candidate configurations
are screened for 120 rounds and shortlisted configurations are confirmed
over 300 rounds using two validation seeds. The selected configuration
uses classifier learning rate 0.1, autoencoder learning rate
$3\times10^{-4}$, and clipping norm 10.

The classifier optimizer is SGD with zero momentum and weight decay.
Encoder and decoder use Adam with $\beta=(0.9,0.999)$,
$\epsilon=10^{-8}$ and zero weight decay. Optimizers are initialized at
each client participation, with Adam state retained between the two
local phases. Learning rates are constant. Computation uses Float32
with AMP and TF32 disabled.

The selected configurations are fixed before the final runs.
Both methods are trained from scratch on all 743 training images per
client for 300 rounds, using paired seeds 42--46.

\subsection{Evaluation and results}

At round 300, each client is evaluated on its domain's full test set.
FusedSpaceFed predictions combine that client's encoder with the final
shared decoder and classifier; FedAvg uses its final shared classifier.
Evaluation uses inference mode and the parameters obtained from training.

Let $c_{s,d}$ and $n_d$ denote the correct-prediction count and test-set
size for seed $s$ and domain $d$, respectively. The uniform-domain
primary accuracy and the sample-weighted accuracy are, respectively,
\[
 U_s=\frac{100}{5}\sum_{d=1}^{5}\frac{c_{s,d}}{n_d},
 \qquad
 A_s=100\frac{\sum_{d=1}^{5}c_{s,d}}{\sum_{d=1}^{5}n_d}.
\]
We report means and sample standard deviations across the five paired
seeds, using the $n-1$ denominator. The uniform-domain accuracy is
$85.86\pm0.85\%$ for FusedSpaceFed and $82.06\pm0.22\%$ for FedAvg.
Table~\ref{tab:capacity_compute_seeds} reports the primary result and
paired differences. The sample-weighted result is $84.51\pm0.75\%$
for FusedSpaceFed and $79.68\pm0.39\%$ for FedAvg.
Table~\ref{tab:capacity_compute_domains} gives the domain-level results.

\begin{table}[H]
\coauthorcolor
\centering
\small
\setlength{\tabcolsep}{4pt}
\caption{Final-round uniform-domain accuracy (\%) and paired difference
(FusedSpaceFed minus FedAvg, in percentage points).
Bold indicates the higher accuracy in each row.}
\label{tab:capacity_compute_seeds}
\begin{tabular}{lrrr}
\toprule
Seed & FusedSpaceFed & FedAvg & Difference \\
\midrule
42 & \textbf{84.67} & 82.27 & $+2.40$ \\
43 & \textbf{86.30} & 82.15 & $+4.15$ \\
44 & \textbf{85.32} & 82.25 & $+3.08$ \\
45 & \textbf{86.76} & 81.78 & $+4.98$ \\
46 & \textbf{86.26} & 81.89 & $+4.37$ \\
\midrule
Mean & \textbf{85.86} & 82.06 & $+3.80$ \\
Sample SD & 0.85 & 0.22 & 1.04 \\
\bottomrule
\end{tabular}
\end{table}

\begin{table}[H]
\coauthorcolor
\centering
\small
\setlength{\tabcolsep}{4pt}
\caption{Final-round domain accuracy (\%): mean and sample standard
deviation across five seeds.
Bold indicates the higher mean accuracy in each row.}
\label{tab:capacity_compute_domains}
\begin{tabular}{lrr}
\toprule
Domain & FusedSpaceFed & FedAvg \\
\midrule
MNIST & $\boldsymbol{96.95\pm0.24}$ & $95.93\pm0.24$ \\
SVHN & $\boldsymbol{67.83\pm2.88}$ & $61.57\pm1.43$ \\
USPS & $\boldsymbol{96.67\pm0.28}$ & $95.94\pm0.25$ \\
SynthDigits & $\boldsymbol{86.32\pm0.50}$ & $81.31\pm0.54$ \\
MNIST-M & $\boldsymbol{81.55\pm0.77}$ & $75.58\pm0.64$ \\
\bottomrule
\end{tabular}
\end{table}
